\chapter{Ramy teoretyczne}

\section{Uczenie ze wzmocnieniem}

Uczenie ze wzmocnieniem (ang. \textit{reinforcement learning}, RL) jest paradygmatem uczenia maszynowego, w którym \textbf{decydent (agent)} obserwuje stan $s \in S$ i wybiera akcję $a \in A$, aby maksymalizować przyszłe nagrody. \textbf{Środowisko} opisane jest przez dynamikę $q(s' \mid s,a)$ i funkcję nagrody $r(s,a)$, które określają konsekwencje działań. \textbf{Polityka} $\phi$ to reguła decyzyjna mapująca stany (lub historie) na rozkłady akcji. Klasyczny cykl RL można opisać jako: obserwacja $\rightarrow$ decyzja $\rightarrow$ nagroda $\rightarrow$ nowy stan.

\subsection{Quasi-hiperboliczne uczenie ze wzmocnieniem}

Wersja quasi-hiperboliczna wzbogaca RL o uprzedzenie teraźniejszości (ang. \textit{present bias}): wagi dyskontowe mają postać $d(0)=1$, $d(t)=\alpha\,\beta^{t}$ dla $t \ge 1$. Parametr $\alpha > 0$ zwiększa względną wartość nagród natychmiastowych, a $\beta \in [0,1)$ zachowuje zależność wykładniczą w przyszłości. Optymalizacja polityki wymaga pogodzenia krótkoterminowych impulsów z celami długoterminowymi. Taki model lepiej odzwierciedla decyzje ludzi, którzy preferują wcześniejsze nagrody.

\section{Definicja modelu}

Model markowskiego procesu decyzyjnego (MDP) z dyskontowaniem quasi-hiperbolicznym jest zdefiniowany jako: $M = (S, A, q, r, \alpha, \beta)$, gdzie:

\begin{itemize}
\item $S$ -- zbiór stanów,
\item $A$ -- zbiór akcji,
\item $q(s' \mid s,a)$ -- prawdopodobieństwo przejścia ze stanu $s$ do stanu $s'$ przy akcji $a$,
\item $r : S \times A \rightarrow \mathbb{R}$ -- funkcja nagrody,
\item $\alpha > 0$ -- parametr wzmacniający wagę nagród natychmiastowych,
\item $\beta \in [0,1)$ -- standardowy współczynnik dyskontowania wykładniczego.
\end{itemize}

Dyskontowanie quasi-hiperboliczne jest zadane przez: $d(0) = 1$, a dla $t \ge 1$ mamy $d(t) = \alpha\,\beta^t$.

\section{Przestrzenie pomocnicze}

Oznaczmy przez:
\begin{enumerate}
\item $\Pr(A)$ -- zbiór miar probabilistycznych na $A$; dla skończonego $A$ utożsamiamy z $\mathbb{R}^{|A|}$.
\item $\mathcal{H}_n$ -- historie długości $n$: $h_n = (s_0, a_0, s_1, a_1, \ldots, s_n)$, przy czym $h_0 = s_0$.
\end{enumerate}

Historia gromadzi pełen kontekst interakcji agent--środowisko do czasu $n$.

\section{Polityki decyzyjne}

\begin{enumerate}
\item Ogólna polityka: $\phi = (\phi_0, \phi_1, \ldots)$, gdzie $\phi_n : \mathcal{H}_n \rightarrow \Pr(A)$ określa rozkład akcji dla każdej historii $h_n \in \mathcal{H}_n$, $n \in \mathbb{N} \cup \{0\} = \mathbb{N}_0$.
\item Polityki markowskie: zależne od stanu i czasu, $\phi_n : S \rightarrow \Pr(A)$.
\item Polityki stacjonarne: zależne tylko od stanu, $\phi_s : S \rightarrow \Pr(A)$; będziemy oznaczać je przez $\phi_s$.
\end{enumerate}

Zbiory polityk:
\begin{itemize}
\item $\Phi$ -- zbiór wszystkich polityk zależnych od historii,
\item $\Phi_m$ -- zbiór polityk markowskich,
\item $\Phi_s$ -- zbiór polityk stacjonarnych.
\end{itemize}

\section{Funkcja wypłaty}

Rozważamy oczekiwaną wypłatę w nieskończonym horyzoncie czasowym z dyskontowaniem quasi-hiperbolicznym. Współczynniki w kolejnych okresach wynoszą 1, $\alpha \beta$, $\alpha \beta^2$, $\alpha \beta^3$, \ldots.

Wypłata dyskontowana w nieskończonym horyzoncie czasowym dla polityki $\phi \in \Phi$ i stanu początkowego $s_0$ jest dana wyrażeniem
\begin{align}
    J_{\phi}^{\alpha,\beta}(s_0)
    &= \mathbb{E}_{\substack{a_0 \sim \phi_0(\cdot \mid s_0) \\ s' \sim q(\cdot \mid s_0,a_0)}}
      \left[ r(s_0,a_0) + \alpha\,\beta\, J_\phi^{\beta}(s') \right] \nonumber \\
    &= \sum_a \phi_0(a \mid s_0) \left[ r(s_0,a) + \alpha\,\beta\, \sum_{s'} q(s' \mid s_0,a) J_\phi^{\beta}(s') \right],
\label{eq:qh-return}
\end{align}
gdzie $\overrightarrow{\phi} = (\phi_1, \phi_2, \ldots)$, a $J_{\phi}^{\beta}(s')$ oznacza standardową wypłatę oczekiwaną z dyskontowaniem wykładniczym
\begin{equation}
J_\phi^{\beta}(s') = \sum_{n=1}^{\infty} \beta^{n-1} \; r^{(n)}(s', \overrightarrow{\phi}).
\label{eq:beta-value}
\end{equation}

\subsection{Oczekiwana nagroda w kroku n}

Funkcja $r^{(n)}(s', \overrightarrow{\phi})$ oznacza oczekiwaną nagrodę w $n$-tym kroku procesu:
\begin{equation}
r^{(n)}(s', \overrightarrow{\phi}) = \mathbb{E}_{\overrightarrow{\phi}}\!\left[ r(s_n, a_n) \;\middle|\; s_0 = s' \right].
\end{equation}
Jest to oczekiwanie po trajektoriach generowanych przez politykę $\overrightarrow{\phi} = (\phi_1, \phi_2, \ldots)$.

Dla $n = 1$:
\begin{align*}
  r^{(1)}(s_1, \overrightarrow{\phi})
  &= \sum_{a_0}\sum_{s_1}\sum_{a_1}\,
  r(s_1, a_1)\,
  \phi_1(a_1 \mid s_1,a_0,s_0)\,
    q(s_1 \mid s_0, a_0)\,
    \phi_0(a_0 \mid s_0).
\end{align*}

Dla $n = 2$:
\begin{align*}
  r^{(2)}(s_2, \overrightarrow{\phi})
    &= \sum_{a_0}\sum_{s_1}\sum_{a_1}\sum_{s_2}\sum_{a_2}
      r(s_2, a_2)\,
      \phi_2(a_2 \mid s_2,a_1,s_1,a_0,s_0) \,
      q(s_2 \mid s_1, a_1)\, \\
      &\qquad{}\phi_1(a_1 \mid s_1,a_0,s_0)\,
      q(s_1 \mid s_0, a_0)\,
      \phi_0(a_0 \mid s_0).
\end{align*}

\section{Struktura optymalnej polityki}

\begin{theorem}[Redukcja polityk]
\label{thm:policy-reduction}
Dla dowolnej $\phi \in \Phi$ istnieją reguły $\mu$ i $\phi_s$ takie, że
\[
  J_{\phi}^{\alpha,\beta}(s) = J_{(\mu,\phi_s)}^{\alpha,\beta}(s), \qquad \forall s \in S.
\]
\end{theorem}

Pierwszy krok steruje dowolna reguła $\mu = \phi_0$, kolejne kroki wykorzystują stacjonarną $\phi_s$. Analiza optymalności sprowadza się do klas polityk dwuskładnikowych $(\mu,\phi_s)$.

\begin{proof}
Dla dowolnej historycznej polityki $\phi = (\phi_0, \phi_1, \ldots)$ mamy
\[
  J_{\phi}^{\alpha,\beta}(s) = \mathbb{E}\!\left[ \sum_{t=0}^{\infty} d(t)\, r(s_t, a_t) \;\middle|\; s_0 = s \right],
\]
gdzie $d(0) = 1$ oraz $d(t) = \alpha\,\beta^{t}$ dla $t \ge 1$.

Rozwijając wartość oczekiwaną otrzymujemy
\[
  J_{\phi}^{\alpha,\beta}(s) = \mathbb{E}\!\left[ r(s_0,a_0) + \sum_{t=1}^{\infty} \alpha \beta^{t}\, r(s_t,a_t) \;\middle|\; s_0 = s \right].
\]

Warunkując po pierwszej akcji $a_0 \sim \phi_0(\cdot \mid s)$ dostajemy
\[
  J_{\phi}^{\alpha,\beta}(s) = \mathbb{E}_{a_0 \sim \phi_0(\cdot \mid s)}\!\left[ r(s,a_0) + \,\mathbb{E}\!\left[ \sum_{t=1}^{\infty} \alpha \beta^{t}\, r(s_t,a_t) \;\middle|\; s_0 = s, a_0 \right] \right].
\]

Kolejne warunkowanie względem $s' \sim q(\cdot \mid s,a_0)$ daje
\[
  J_{\phi}^{\alpha,\beta}(s) = \mathbb{E}_{\substack{a_0 \sim \phi_0(\cdot \mid s) \\ s' \sim q(\cdot \mid s,a_0)}}\!\left[ r(s,a_0) + \alpha\beta\,\mathbb{E}\!\left[ \sum_{t=1}^{\infty} \beta^{t-1}\, r(s_t,a_t) \;\middle|\; s_1=s' \right] \right].
\]

Pozostała suma dyskontowana ma postać wykładniczą, więc
\[
  \mathbb{E}\!\left[ \sum_{t=1}^{\infty} \beta^{t-1}\, r(s_t,a_t) \;\middle|\; s_1 = s' \right] = \,J_{\overrightarrow{\phi}}^{\beta}(s_1),
\]
gdzie $\overrightarrow{\phi} = (\phi_1, \phi_2, \ldots)$ i $J^{\beta}$ odpowiada klasycznemu dyskontowaniu wykładniczemu.

Dla takiego dyskontowania istnieje stacjonarna polityka $\phi_s \in \Phi_s$, która spełnia
\[
  J^{\beta}_{\overrightarrow{\phi}}(s') = J^{\beta}_{\phi_s}(s'), \qquad \forall s' \in S.
\]

Podstawiając do poprzedniego równania, otrzymujemy
\[
  J_{\phi}^{\alpha,\beta}(s) = \sum_{a}\phi_{0}(a \mid s)\!\left[ r(s,a) + \alpha\,\beta\, \sum_{s'} q(s' \mid s,a) J^{\beta}_{\phi_s}(s') \right].
\]

Wybierając $\mu = \phi_0$ jako regułę pierwszego kroku, uzyskujemy żądaną równość
\[
  J_{\phi}^{\alpha,\beta}(s) = J_{(\mu,\phi_s)}^{\alpha,\beta}(s), 
\]
co kończy dowód.
\end{proof}

\section{Optymalna polityka}

W modelu z dyskontowaniem quasi-hiperbolicznym optymalne polityki mają specyficzną strukturę, różniącą się od przypadku wykładniczego. 

\subsection{Definicja optymalnej polityki}

Optymalna para $(\mu^\star, \phi_s^\star) \in \Phi_m \times \Phi_s$ maksymalizuje wartość w każdym stanie:
\begin{equation}
V^{\alpha, \beta}_{\mu^\star, \phi_s^\star}(s) = \max_{\mu \in \Phi_m} \max_{\phi_s \in \Phi_s} V^{\alpha, \beta}_{\mu, \phi_s}(s), \quad \forall s \in S.
\end{equation}

\subsection{Wyprowadzenie struktury optymalnej polityki}

Punktem wyjścia jest redukcja polityk (Twierdzenie~\ref{thm:policy-reduction}):
\begin{align}
V^{\alpha, \beta}_{\mu^\star, \phi_s^\star}(s) &= \max_{\mu \in \Phi_m} \max_{\phi_s \in \Phi_s} \sum_{a \in A} \mu(a\mid s) \Big[ r(s, a) + \alpha \beta \sum_{s' \in S} q(s'\mid s, a) V^{\beta}_{\phi_s}(s') \Big] \nonumber \\
&= \max_{\mu \in \Phi_m} \sum_{a \in A} \mu(a\mid s) \Big[ r(s, a) + \alpha \beta \sum_{s'} q(s'\mid s, a) \max_{\phi_s \in \Phi_s} V^{\beta}_{\phi_s}(s') \Big].
\end{align}

Dla dyskontowania wykładniczego optymalna jest $\pi_\beta^\star$, stąd $\phi_s^\star = \pi_\beta^\star$ oraz $V^{\beta}_{\phi_s^\star} = V^{\beta}_\star$.

Optymalna reguła pierwszego kroku:
\begin{equation}
\mu^\star(s) = \arg\max_{a \in A} \left[ r(s, a) + \alpha \beta \sum_{s'} q(s'\mid s, a) V^{\beta}_{\pi^\star_\beta}(s') \right].
\end{equation}

\begin{theorem}[Deterministyczna optymalna polityka]
\label{thm:deterministic-optimal}
Istnieje deterministyczna optymalna para $(\mu^\star, \phi_s^\star)$ taka, że
\begin{equation}
V^{\alpha, \beta}_{\mu^\star, \phi_s^\star}(s) = \max_{\mu \in \Phi_m} \max_{\phi_s \in \Phi_s} V^{\alpha, \beta}_{\mu, \phi_s}(s), \quad \forall s \in S.
\end{equation}
\end{theorem}

\begin{proof}
Polityka $\phi_s^\star$ jest deterministyczna, ponieważ $\pi_\beta^\star$ w klasycznym MDP taką jest. Optymalizacja $\mu$ nad sympleksem $\Pr(A)$ jest liniowa, więc optimum znajduje się w wierzchołku, co oznacza deterministyczne $\mu^\star$.
\end{proof}

W praktyce wystarczy przeszukiwać $\mu^\star \in A^S$ oraz $\phi_s^\star \in A^S$.